% Auto-generated complete chapter from 23_interview_answer_bank.md
\chapter{23\_interview\_answer\_bank}
\label{complete:root_027}

\begin{sourcemeta}
\textbf{Fuente:} \href{file:///E:/Laboral/23_interview_answer_bank.md}{23\_interview\_answer\_bank.md}\\
\textbf{Seccion:} root\\
\textbf{Estado:} canonical\\
\textbf{Rol:} Modulo principal\\
\textbf{Lineas originales:} 902\\
\textbf{Ultima modificacion:} 2026-06-11 15:25\\
\textbf{Deduplicacion conservadora:} 0 bloques repetidos omitidos; 0 caracteres repetidos no reimpresos.
\end{sourcemeta}

\section*{Resumen de orientacion}
This document provides prepared interview answers for Alexander Woodcock Salomón’s target roles.

\section*{Contenido completo depurado}
\addcontentsline{toc}{section}{Contenido completo depurado}




\subsection{Purpose}




This document provides prepared interview answers for Alexander Woodcock Salomón’s target roles.




It converts the strategy, CV, LinkedIn, portfolio, TwinSight case study, negotiation framework and outreach documents into interview-ready language.




The answers are written primarily in English because the target market is international. Spanish versions are included where they are useful for LATAM or Spanish-speaking recruiters.




\medskip\hrule\medskip




\section{1. Interview positioning rule}




Every interview answer should reinforce the same profile:




\begin{mdcode}
Bloque de codigo: text
Real-Time 3D Developer / Unity Technical Artist focused on interactive technical visualization, CAD-to-realtime optimization, and Unity WebGL deployment.
\end{mdcode}




Do not drift into:




\begin{itemize}
\item generic full-stack developer;
\item pure 3D artist;
\item AI engineer;
\item game designer;
\item graphic designer;
\item general multimedia generalist.
\end{itemize}




\medskip\hrule\medskip




\section{2. Core interview narrative}




\subsection{30-second version}




\begin{mdcode}
Bloque de codigo: text
I’m a Colombia-based Real-Time 3D Developer / Unity Technical Artist focused on Unity WebGL, C#, technical visualization and CAD-to-realtime optimization. My strongest project is TwinSight X500, a Unity WebGL drone assembly inspection prototype with component selection, exploded view, cross-section tools, visual modes, technical UI and usability/workload evaluation. I’m now looking for remote roles where I can apply that background to interactive 3D, technical visualization, simulation, digital twin, XR-adjacent or Unity technical art work.
\end{mdcode}




\subsection{60-second version}




\begin{mdcode}
Bloque de codigo: text
I’m a Real-Time 3D Developer / Unity Technical Artist based in Colombia. My focus is building interactive technical visualization systems with Unity, C#, WebGL and Blender.

My flagship project is TwinSight X500, a Unity WebGL prototype for inspecting the assembly of a Holybro X500 V2 drone. The project converts technical documentation and CAD-derived assets into an optimized real-time 3D experience with component selection, exploded view, cross-section tools, visual modes and technical information panels.

My background combines Multimedia Engineering, advanced Electronic Engineering coursework, real-time 3D, Blender-based asset preparation, technical documentation and Python-assisted tooling. I’m looking for remote contractor or full-time roles where Unity, interactive 3D, technical visualization, WebGL or simulation/digital twin interfaces are relevant.
\end{mdcode}




\subsection{90-second version}




\begin{mdcode}
Bloque de codigo: text
I’m a Colombia-based Real-Time 3D Developer / Unity Technical Artist focused on interactive technical visualization, CAD-to-realtime optimization and Unity WebGL deployment.

My strongest project is TwinSight X500, a Multimedia Engineering thesis project built in Unity WebGL to inspect and understand the assembly of a Holybro X500 V2 drone. The project addresses a common issue in technical documentation: 2D diagrams can show components, but users often have to mentally reconstruct spatial relationships.

TwinSight turns that into an interactive 3D viewer. Users can select components, inspect technical information, use exploded views, apply cross-section tools and switch between visual modes. A major part of the project was converting CAD-derived geometry into optimized real-time assets using Blender-based cleanup, preparation and Unity integration.

I also have secondary Python tooling work through ARA Framework, a research automation prototype. I treat that as supporting evidence of automation and systems thinking, while the core direction remains Unity, real-time 3D and technical visualization.

I’m not positioning myself as a senior industry technical artist. I’m positioning myself as a junior-strong / early-mid candidate with unusually relevant project evidence in Unity WebGL and technical visualization.
\end{mdcode}




\medskip\hrule\medskip




\section{3. “Tell me about yourself”}




\subsection{English answer}




\begin{mdcode}
Bloque de codigo: text
I’m a Real-Time 3D Developer / Unity Technical Artist based in Colombia, focused on Unity WebGL, C#, technical visualization and CAD-to-realtime optimization.

My strongest project is TwinSight X500, a Unity WebGL technical visualization prototype for drone assembly inspection. It includes component selection, exploded view, cross-section tools, visual modes, technical UI and an optimized CAD-to-realtime asset workflow.

My academic background is in Multimedia Engineering, with advanced previous coursework in Electronic Engineering, so I tend to approach 3D work from both a visual and engineering perspective. I’m especially interested in roles where real-time 3D is used to make technical systems easier to inspect, understand and communicate.
\end{mdcode}




\subsection{Spanish answer}




\begin{mdcode}
Bloque de codigo: text
Soy Real-Time 3D Developer / Unity Technical Artist en Colombia, enfocado en Unity WebGL, C#, visualización técnica y optimización CAD-to-realtime.

Mi proyecto más fuerte es TwinSight X500, un prototipo Unity WebGL de visualización técnica para inspección de ensamblaje de un dron. Incluye selección de componentes, vista explotada, herramientas de corte, modos visuales, UI técnica y un flujo optimizado de assets derivados de CAD.

Mi formación es en Ingeniería Multimedia, con estudios avanzados previos en Ingeniería Electrónica, así que tiendo a abordar el 3D desde una perspectiva visual y también ingenieril. Me interesan especialmente roles donde el 3D en tiempo real ayude a inspeccionar, entender y comunicar sistemas técnicos.
\end{mdcode}




\subsection{What not to say}




\begin{mdcode}
Bloque de codigo: text
I’m a multimedia engineer who does a bit of everything: 3D, AI, web, design, games, fitness apps and research.
\end{mdcode}




This dilutes the positioning.




\medskip\hrule\medskip




\section{4. “Explain TwinSight X500”}




\subsection{30-second answer}




\begin{mdcode}
Bloque de codigo: text
TwinSight X500 is a Unity WebGL technical visualization prototype for inspecting the assembly of a Holybro X500 V2 drone. It transforms technical documentation and CAD-derived assets into an interactive 3D viewer with component selection, exploded view, cross-section tools, visual modes and technical information panels. The goal is to help users understand spatial relationships in the assembly, not just look at individual parts.
\end{mdcode}




\subsection{60-second answer}




\begin{mdcode}
Bloque de codigo: text
TwinSight X500 is my Multimedia Engineering thesis project. It is a Unity WebGL interactive 3D visualization prototype for inspecting and understanding the assembly of a Holybro X500 V2 drone.

The problem is that traditional assembly documentation is usually static. It can show parts and diagrams, but the user still has to mentally reconstruct how components relate in space.

TwinSight addresses that through an interactive real-time 3D viewer. Users can select components, inspect technical information, use exploded view to understand relationships, apply cross-section tools and switch between visual modes.

Technically, it involved Unity/C# runtime systems, UI Toolkit, WebGL deployment, Blender-based asset preparation and CAD-to-realtime optimization.
\end{mdcode}




\subsection{Technical answer}




\begin{mdcode}
Bloque de codigo: text
TwinSight X500 is a browser-accessible Unity WebGL prototype focused on technical inspection of a drone assembly.

The system combines a CAD-to-realtime asset pipeline with runtime interaction systems. On the asset side, dense CAD-derived geometry had to be cleaned, simplified and prepared for real-time use. On the Unity side, the prototype includes component selection, highlighting, technical information panels, exploded view behavior, cross-section/clipping tools and multiple visual modes.

The project was also evaluated academically using usability and workload methods such as SUS, NASA-TLX Raw and Think-Aloud. The important point is that the project is not only a model viewer; it is an interaction system for technical understanding.
\end{mdcode}




\medskip\hrule\medskip




\section{5. “What problem does TwinSight solve?”}




\subsection{Answer}




\begin{mdcode}
Bloque de codigo: text
The project addresses the gap between static technical documentation and spatial understanding.

In assembly documentation, users can see parts in 2D diagrams or manuals, but they often have to mentally reconstruct how components relate in 3D space. That creates cognitive load, especially when there are overlapping parts, internal structures, fasteners or subassemblies.

TwinSight explores whether an interactive 3D viewer can reduce that burden by allowing users to directly inspect the assembly, select components, see exploded relationships, use cross-sections and switch visual modes.
\end{mdcode}




\subsection{Short version}




\begin{mdcode}
Bloque de codigo: text
It helps users understand spatial relationships in a technical assembly instead of relying only on static 2D diagrams.
\end{mdcode}




\medskip\hrule\medskip




\section{6. “What did you personally build?”}




\subsection{Answer}




\begin{mdcode}
Bloque de codigo: text
My work covered Unity/C# implementation, runtime interaction systems, technical UI, asset integration, Blender-based asset preparation, CAD-to-realtime optimization, documentation, evaluation design and final integration.

Specifically, I worked on systems such as component selection, exploded view behavior, cross-section/clipping tools, visual modes and technical information panels.

I also handled the process of converting CAD-derived assets into something more appropriate for real-time Unity WebGL deployment. AI tools were used as support during parts of implementation and debugging, but I remained responsible for architecture decisions, integration, validation, technical judgment and final ownership.
\end{mdcode}




\subsection{If challenged}




\begin{mdcode}
Bloque de codigo: text
I can walk through the project architecture, explain the feature systems, describe the asset pipeline and discuss what I would improve. My claim is not that I wrote every line without assistance; my claim is that I owned the technical integration, decisions and delivery of the system.
\end{mdcode}




\medskip\hrule\medskip




\section{7. “How much of this did you build yourself?”}




\subsection{Answer}




\begin{mdcode}
Bloque de codigo: text
I was responsible for the project direction, architecture decisions, Unity integration, feature implementation, debugging, asset workflow decisions, documentation, evaluation and final delivery.

I did use AI-assisted tools during parts of the process, mainly for implementation support, debugging, research and documentation. I treat that similarly to using a technical assistant: useful for acceleration, but not a replacement for understanding the system.

I can explain the project structure, the technical decisions, the tradeoffs, the limitations and the areas I would refactor.
\end{mdcode}




\subsection{Short version}




\begin{mdcode}
Bloque de codigo: text
I used AI as a support tool, but I owned the architecture, integration, validation and final technical decisions.
\end{mdcode}




\subsection{What not to say}




\begin{mdcode}
Bloque de codigo: text
AI did most of it.
\end{mdcode}




or:




\begin{mdcode}
Bloque de codigo: text
I wrote every line manually.
\end{mdcode}




Both are risky if inaccurate.




\medskip\hrule\medskip




\section{8. “How did you use AI?”}




\subsection{Answer}




\begin{mdcode}
Bloque de codigo: text
I used AI-assisted tools mainly for implementation support, debugging, research, documentation drafting and exploring solution alternatives.

I did not use AI as a substitute for technical ownership. I was still responsible for choosing the architecture, integrating the systems, testing behavior, validating outputs and deciding what was technically appropriate for the project.

The best way to describe it is that AI helped me move faster, but I still had to understand, adapt, debug and own the result.
\end{mdcode}




\subsection{If interviewer is skeptical}




\begin{mdcode}
Bloque de codigo: text
That concern is valid. That is why I do not present the project as proof that I manually authored every line in isolation. I present it as proof that I can define a technical problem, integrate Unity systems, prepare assets, validate behavior, document tradeoffs and use modern tools responsibly to deliver a working prototype.
\end{mdcode}




\medskip\hrule\medskip




\section{9. “Why this role?”}




\subsection{Unity / Real-Time 3D answer}




\begin{mdcode}
Bloque de codigo: text
This role fits because it combines Unity development, real-time interaction and technical problem solving. My strongest project, TwinSight X500, required Unity/C# runtime systems, WebGL deployment, technical UI, asset optimization and real-time inspection features.

I am specifically looking for roles where Unity is used to build interactive systems, not just passive visuals.
\end{mdcode}




\subsection{Technical Visualization answer}




\begin{mdcode}
Bloque de codigo: text
This role fits because my strongest project is directly about turning technical information into an interactive 3D visualization system. TwinSight converts drone assembly documentation and CAD-derived assets into a real-time inspection experience with selection, exploded view, cross-section and visual modes.

That aligns with technical visualization, digital twins, simulation interfaces and product/assembly visualization.
\end{mdcode}




\subsection{Technical Artist answer}




\begin{mdcode}
Bloque de codigo: text
This role fits because my work sits between 3D asset workflows and runtime systems. In TwinSight I had to prepare and optimize 3D assets, integrate them into Unity, build inspection tools, create visual modes and consider WebGL constraints.

I’m not positioning myself as a senior shader specialist, but as a Unity technical artist profile focused on optimization, runtime interaction, technical UI and visualization workflows.
\end{mdcode}




\subsection{WebGL answer}




\begin{mdcode}
Bloque de codigo: text
This role fits because TwinSight was built for Unity WebGL deployment. The project required thinking about browser delivery, asset weight, interaction design and performance constraints.

I’m interested in building browser-accessible 3D tools where users can interact with technical or product data without needing native CAD or desktop tools.
\end{mdcode}




\medskip\hrule\medskip




\section{10. “Why should we hire you?”}




\subsection{Main answer}




\begin{mdcode}
Bloque de codigo: text
You should consider me if the role needs someone who can bridge Unity development, real-time 3D asset workflows and technical visualization.

My strongest evidence is TwinSight X500, where I built a Unity WebGL technical visualization prototype for drone assembly inspection. It required C# interaction systems, WebGL-oriented optimization, technical UI, CAD-to-realtime asset preparation, documentation and usability/workload evaluation.

I’m not claiming senior industry experience. The value I bring is a focused technical profile, strong project evidence, English communication, engineering-oriented thinking and a clear interest in building real-time 3D systems for technical use cases.
\end{mdcode}




\subsection{More concise version}




\begin{mdcode}
Bloque de codigo: text
Because I have direct project evidence for the type of work this role needs: Unity, WebGL, real-time 3D interaction, technical visualization, asset optimization and documentation. I’m not a generic junior developer; I have a focused technical visualization project that is highly aligned with this kind of role.
\end{mdcode}




\medskip\hrule\medskip




\section{11. “Why no formal industry experience in this field yet?”}




\subsection{Answer}




\begin{mdcode}
Bloque de codigo: text
My formal industry experience in technical art and Unity is still limited because my recent focus has been completing my Multimedia Engineering thesis and building a portfolio around real-time 3D and technical visualization.

However, the thesis project is substantial and directly relevant. TwinSight X500 required Unity/C# implementation, WebGL deployment, CAD-to-realtime optimization, runtime interaction systems, technical UI, documentation and evaluation.

So I’m not positioning myself as a senior industry candidate. I’m positioning myself as a junior-strong / early-mid candidate with unusually relevant project evidence in Unity technical visualization.
\end{mdcode}




\subsection{Spanish version}




\begin{mdcode}
Bloque de codigo: text
Mi experiencia formal de industria en technical art y Unity todavía es limitada porque mi foco reciente fue completar la tesis de Ingeniería Multimedia y construir un portafolio alrededor de 3D en tiempo real y visualización técnica.

Sin embargo, el proyecto de tesis es sustancial y directamente relevante. TwinSight X500 requirió implementación en Unity/C#, despliegue WebGL, optimización CAD-to-realtime, sistemas de interacción runtime, UI técnica, documentación y evaluación.

No me presento como candidato senior. Me posiciono como un perfil junior fuerte / early-mid con evidencia de proyecto muy relevante en visualización técnica con Unity.
\end{mdcode}




\medskip\hrule\medskip




\section{12. “What are your strongest technical skills?”}




\subsection{Answer}




\begin{mdcode}
Bloque de codigo: text
My strongest technical areas are Unity/C# runtime interaction systems, Unity WebGL deployment, technical visualization, CAD-to-realtime asset optimization and Blender-based asset preparation.

In practical terms, that means I can work on interactive 3D inspection tools, component selection systems, technical UI, exploded views, cross-section tools, visual modes and optimization workflows for real-time deployment.
\end{mdcode}




\subsection{Add if role is tools/Python}




\begin{mdcode}
Bloque de codigo: text
I also have secondary experience with Python automation and LLM workflow prototyping through ARA Framework, which supports tooling and documentation workflows.
\end{mdcode}




\medskip\hrule\medskip




\section{13. “What are your weaknesses or gaps?”}




\subsection{Safe answer}




\begin{mdcode}
Bloque de codigo: text
My main gap is formal industry production experience in a professional Unity or technical art team. Most of my strongest evidence comes from thesis and portfolio work rather than shipped commercial products.

I’m addressing that by making my work public, documenting it clearly, preparing technical breakdowns and targeting roles where my project evidence is directly relevant.

Another gap is that I am not yet a deep shader/HLSL specialist. I have Shader Graph and basic HLSL familiarity, but I position that as a growth area rather than an expert claim.
\end{mdcode}




\subsection{Alternative}




\begin{mdcode}
Bloque de codigo: text
I still need more professional team experience, code review exposure and production-scale collaboration. My current strength is project ownership and technical integration. The next step is applying that in a team environment with production standards.
\end{mdcode}




\medskip\hrule\medskip




\section{14. “What would you improve in TwinSight?”}




\subsection{Answer}




\begin{mdcode}
Bloque de codigo: text
I would improve five areas.

First, I would do deeper profiling and optimization, especially for WebGL build size, memory use and frame stability.

Second, I would separate the runtime systems more cleanly into modules, especially selection, exploded view, clipping and visual modes.

Third, I would improve the public WebGL deployment pipeline and add clearer build documentation.

Fourth, I would expand the evaluation with more participants and more controlled task design.

Fifth, I would refine the asset pipeline with more formal LODs and stricter asset validation.
\end{mdcode}




\subsection{Short version}




\begin{mdcode}
Bloque de codigo: text
I would improve profiling, modular architecture, WebGL build pipeline, evaluation depth and LOD/asset validation.
\end{mdcode}




\medskip\hrule\medskip




\section{15. “What was the hardest technical problem?”}




\subsection{Answer}




\begin{mdcode}
Bloque de codigo: text
The hardest problem was balancing technical asset fidelity with real-time WebGL constraints.

CAD-derived geometry can be extremely dense and not appropriate for real-time browser deployment. The challenge was to preserve component identity and assembly relationships while reducing complexity enough for interactive performance.

That required asset inspection, cleanup, simplification, Blender-based preparation, Unity integration and constant tradeoffs between visual clarity, performance and scope.
\end{mdcode}




\subsection{Alternative: interaction-system answer}




\begin{mdcode}
Bloque de codigo: text
Another difficult part was designing the runtime interaction model so that the viewer was not just a 3D model on screen. The features had to support inspection: component selection, highlighting, information panels, exploded view, cross-section and visual modes all needed to work as part of a coherent technical visualization workflow.
\end{mdcode}




\medskip\hrule\medskip




\section{16. “How did you approach optimization?”}




\subsection{Answer}




\begin{mdcode}
Bloque de codigo: text
I approached optimization from the asset pipeline first.

The source geometry was CAD-derived, so the initial priority was reducing unnecessary geometric density while preserving component identity and assembly readability. That involved cleanup, simplification, retopology decisions, UV preparation and Unity import decisions.

On the Unity side, I considered WebGL constraints, visual clarity, material complexity, interaction responsiveness and scene organization.

The goal was not just fewer polygons. The goal was a model that remained technically understandable while being suitable for browser-based real-time interaction.
\end{mdcode}




\subsection{If asked for metrics}




\begin{mdcode}
Bloque de codigo: text
The reported reduction was from over 6.5 million triangles to approximately 95,617 optimized triangles, but I would verify the final number against the thesis report before publishing it as a fixed public claim.
\end{mdcode}




\medskip\hrule\medskip




\section{17. “How would you debug low FPS in Unity WebGL?”}




\subsection{Answer}




\begin{mdcode}
Bloque de codigo: text
I would start by separating whether the bottleneck is CPU, GPU, memory, loading or browser-specific overhead.

First, I would profile in Unity and then test the WebGL build in the browser because editor performance can be misleading. I would check draw calls, triangle count, material count, texture sizes, overdraw, scripts running per frame and garbage allocation.

For WebGL specifically, I would also check build size, memory configuration, compression, asset loading, shader complexity and browser console warnings.

Then I would apply fixes based on evidence: reduce geometry, combine or simplify materials, optimize textures, remove unnecessary update loops, avoid allocations, simplify shaders, use LODs where appropriate and reduce UI or rendering overhead.
\end{mdcode}




\medskip\hrule\medskip




\section{18. “How would you structure an exploded view system?”}




\subsection{Answer}




\begin{mdcode}
Bloque de codigo: text
I would treat exploded view as a controlled transformation system.

Each component or group needs an original position and an exploded target position. The system should store those transforms, interpolate between them, and preserve assembly relationships through grouping or hierarchy.

For a technical visualization project, I would avoid random outward movement. The exploded direction should communicate structure, so it should be based on assembly logic, axes, component groups or manually authored offsets.

I would also keep it reversible, expose a slider or state control, and ensure selection/highlighting still works in both assembled and exploded states.
\end{mdcode}




\medskip\hrule\medskip




\section{19. “How would you implement component selection?”}




\subsection{Answer}




\begin{mdcode}
Bloque de codigo: text
I would use colliders, raycasting and a mapping between scene objects and component metadata.

The selection system should not only identify a mesh. It should connect the selected object to a data model containing name, category, technical information, hierarchy, parent/child relationships and UI behavior.

I would also separate visual feedback from selection logic. That way highlighting, isolation, ghosting or UI updates can be changed without rewriting the selection core.
\end{mdcode}




\medskip\hrule\medskip




\section{20. “How would you implement cross-section / clipping?”}




\subsection{Answer}




\begin{mdcode}
Bloque de codigo: text
At a high level, I would use a clipping plane or shader-based clipping approach, depending on the rendering pipeline and platform constraints.

The system needs a plane position, orientation and user control. Materials or shaders then discard fragments on one side of the plane. For Unity WebGL, I would keep the shader approach as simple as possible and test compatibility carefully.

If the project requires capped cross-sections or engineering-grade sectioning, that becomes more complex and may require mesh processing or specialized rendering techniques. For a visualization prototype, shader-based clipping is often a practical starting point.
\end{mdcode}




\medskip\hrule\medskip




\section{21. “What is CAD-to-realtime?”}




\subsection{Answer}




\begin{mdcode}
Bloque de codigo: text
CAD-to-realtime is the process of converting CAD or engineering-derived geometry into assets that can run efficiently in a real-time engine like Unity.

CAD models are usually built for precision and manufacturing, not for real-time rendering. They can contain dense geometry, excessive detail, problematic topology, many small parts, or structures that are inefficient for WebGL or real-time engines.

The workflow usually involves conversion, cleanup, simplification, retopology where needed, UV preparation, material planning, hierarchy organization, optimization and engine integration.

The goal is to preserve the technical meaning of the object while making it usable in an interactive real-time environment.
\end{mdcode}




\medskip\hrule\medskip




\section{22. “What is the difference between technical visualization and a 3D model viewer?”}




\subsection{Answer}




\begin{mdcode}
Bloque de codigo: text
A model viewer primarily lets users look at a 3D object.

Technical visualization should help users understand technical information. That means the interaction, UI, visual modes and data presentation are designed around inspection, comparison, assembly logic or decision-making.

In TwinSight, the goal was not just to display a drone. The goal was to support understanding through component selection, exploded view, cross-section tools, technical information panels and evaluation of usability and workload.
\end{mdcode}




\medskip\hrule\medskip




\section{23. “What is your experience with WebGL?”}




\subsection{Answer}




\begin{mdcode}
Bloque de codigo: text
My direct WebGL experience is through Unity WebGL deployment in TwinSight X500.

The project required thinking about browser constraints, asset optimization, UI readability, interaction performance and deployment limitations. I would not claim to be a low-level WebGL graphics programmer, but I understand the practical constraints of delivering real-time 3D through the browser using Unity WebGL.

I’m interested in expanding that into Three.js/WebGPU where relevant.
\end{mdcode}




\medskip\hrule\medskip




\section{24. “What is your experience with XR?”}




\subsection{Answer}




\begin{mdcode}
Bloque de codigo: text
My strongest direct project is not an XR project, but it is adjacent to XR because it involves real-time 3D interaction, technical visualization, spatial understanding, optimized assets and inspection workflows.

For XR roles, I would be strongest where the role needs Unity, interactive 3D, technical UI, simulation interfaces or training/inspection workflows. I would still need to deepen platform-specific XR experience such as OpenXR, device interaction patterns and performance constraints for headsets.
\end{mdcode}




\medskip\hrule\medskip




\section{25. “What is your experience with shaders?”}




\subsection{Answer}




\begin{mdcode}
Bloque de codigo: text
I have experience with Shader Graph and basic HLSL concepts, especially in the context of visual modes and real-time visualization.

I would not position myself as a senior shader specialist yet. My current shader strength is practical integration: using shaders or material modes to support inspection, ghosting, wireframe-like views, blueprint-style visualization and similar technical visual modes.

Deep shader optimization and advanced HLSL are growth areas.
\end{mdcode}




\medskip\hrule\medskip




\section{26. “What is ARA Framework?”}




\subsection{Answer}




\begin{mdcode}
Bloque de codigo: text
ARA Framework is a secondary Python research automation prototype.

It explores multi-agent style workflows for niche analysis, literature research, technical planning, implementation guidance and structured report generation. It uses LangGraph/LangChain-style orchestration concepts, research APIs, document-processing ideas and Markdown output.

I do not present it as a production AI platform. I use it to demonstrate tooling, automation thinking, Python architecture and technical documentation.
\end{mdcode}




\medskip\hrule\medskip




\section{27. “Are you an AI engineer?”}




\subsection{Answer}




\begin{mdcode}
Bloque de codigo: text
No. I would not position myself as an AI engineer in the machine learning or model-training sense.

I have experience using AI-assisted development workflows and building LLM workflow prototypes like ARA Framework. That gives me practical experience with automation, orchestration and technical tooling, but my primary professional direction is still Unity, real-time 3D and technical visualization.

AI tooling is a secondary differentiator, not the main role identity.
\end{mdcode}




\medskip\hrule\medskip




\section{28. “What salary are you expecting?”}




\subsection{Early-stage answer}




\begin{mdcode}
Bloque de codigo: text
At this stage, I would prefer to understand the scope, responsibilities, expected availability, contract structure and seniority level before giving a final number.

For remote contractor work, I usually evaluate compensation based on monthly retainer or hourly rate, depending on scope and time commitment.
\end{mdcode}




\subsection{With range}




\begin{mdcode}
Bloque de codigo: text
For remote contractor work, depending on scope and expected availability, I’m generally targeting the USD [range] per month range, or an equivalent hourly rate.

I’m open to discussing structure if the role is a strong fit in Unity, real-time 3D, WebGL or technical visualization.
\end{mdcode}




\subsection{If pressed for minimum}




\begin{mdcode}
Bloque de codigo: text
For a full-time remote contractor arrangement, my minimum depends on scope and stability, but I’m primarily targeting roles that are meaningfully above local Colombian compensation and aligned with international Unity/technical visualization rates.
\end{mdcode}




\subsection{What not to say}




\begin{mdcode}
Bloque de codigo: text
I can accept anything.
\end{mdcode}




or:




\begin{mdcode}
Bloque de codigo: text
My minimum is USD 1,500.
\end{mdcode}




Do not reveal the floor early unless necessary.




\medskip\hrule\medskip




\section{29. “Can you work remotely from Colombia?”}




\subsection{Answer}




\begin{mdcode}
Bloque de codigo: text
Yes. I’m currently based in Colombia and available for remote contractor/B2B work.

I can work with reasonable overlap with US time zones, and partial overlap with European time zones depending on the schedule. For employee-style arrangements, the company would need to confirm whether it supports EOR or local compliance. For contractor work, I can operate through invoice-based arrangements subject to Colombian tax and administrative requirements.
\end{mdcode}




\medskip\hrule\medskip




\section{30. “Are you authorized to work in the EU?”}




\subsection{Current answer}




\begin{mdcode}
Bloque de codigo: text
Not currently. I am currently based in Colombia and should be treated as available for remote contractor/B2B work.

I expect a Portuguese passport in the medium term, which may support future EU work authorization, but I do not currently represent myself as EU-authorized until that status is legally issued.
\end{mdcode}




\subsection{After Portuguese passport}




Only after it is legally true:




\begin{mdcode}
Bloque de codigo: text
Yes. I hold Portuguese citizenship and can work in the EU/EEA, subject to local registration and tax requirements.
\end{mdcode}




\medskip\hrule\medskip




\section{31. “Are you open to relocation?”}




\subsection{Answer}




\begin{mdcode}
Bloque de codigo: text
My immediate priority is remote work from Colombia, especially contractor/B2B roles.

I am open to relocation if the opportunity is strong enough and the legal structure is clear. Europe is strategically relevant for me in the medium term due to an expected Portuguese passport and a possible Germany route through family residence, but I do not currently claim EU work authorization.
\end{mdcode}




\medskip\hrule\medskip




\section{32. “Can you work US hours?”}




\subsection{Answer}




\begin{mdcode}
Bloque de codigo: text
Yes, I can work with reasonable overlap with US time zones. I’m in the America/Bogota time zone, which aligns well with Eastern and Central time and can partially overlap with Pacific time depending on the schedule.
\end{mdcode}




\medskip\hrule\medskip




\section{33. “Can you work European hours?”}




\subsection{Answer}




\begin{mdcode}
Bloque de codigo: text
I can support partial overlap with European time zones from Colombia, especially for meetings scheduled earlier in the European day. For a fully Europe-aligned schedule, I would need to review the expectations and sustainability.
\end{mdcode}




\medskip\hrule\medskip




\section{34. “What type of contract are you open to?”}




\subsection{Answer}




\begin{mdcode}
Bloque de codigo: text
I’m open to remote contractor/B2B arrangements, full-time roles where the company can support the appropriate legal structure, and project-based contracts if the scope is clear.

For international work from Colombia, contractor/B2B is usually the most practical short-term structure.
\end{mdcode}




\medskip\hrule\medskip




\section{35. “How do you handle contractor work?”}




\subsection{Answer}




\begin{mdcode}
Bloque de codigo: text
For contractor work, I prefer a clear agreement covering scope, deliverables, expected availability, payment terms, IP ownership, confidentiality, communication cadence and termination terms.

I also separate contractor arrangements from employee-like arrangements. If the company expects fixed full-time hours, exclusivity and direct supervision, then the legal structure needs to be reviewed carefully.
\end{mdcode}




\medskip\hrule\medskip




\section{36. “What are your long-term goals?”}




\subsection{Answer}




\begin{mdcode}
Bloque de codigo: text
My long-term goal is to specialize in real-time 3D systems for technical visualization, simulation, digital twins, XR-adjacent workflows and Unity technical art.

In the next 12–24 months, I want to strengthen my professional experience, improve my portfolio with production-level work, deepen Unity optimization and WebGL/XR knowledge, and position myself for higher-value international roles.
\end{mdcode}




\medskip\hrule\medskip




\section{37. “Why not focus on games if you like videogames?”}




\subsection{Answer}




\begin{mdcode}
Bloque de codigo: text
I am interested in games, but I’m being strategic about employability and compensation.

My strongest current evidence, TwinSight X500, is more aligned with technical visualization, simulation, digital twins, industrial/product visualization and Unity-based interactive systems than with pure game production.

Game technical art is still relevant, especially for Unity runtime, tools, shaders and optimization roles, but I do not want to force the portfolio into a pure game direction if the stronger market fit is technical visualization.
\end{mdcode}




\medskip\hrule\medskip




\section{38. “What kind of team environment do you prefer?”}




\subsection{Answer}




\begin{mdcode}
Bloque de codigo: text
I work best in teams where technical expectations are explicit, feedback is direct and work is organized around clear deliverables.

Because my background combines self-directed project work and technical documentation, I value written communication, task clarity, code/project organization and technical review. For remote work, I think clarity is more important than constant meetings.
\end{mdcode}




\medskip\hrule\medskip




\section{39. “Describe a time you solved a difficult problem”}




\subsection{Answer using TwinSight}




\begin{mdcode}
Bloque de codigo: text
A difficult problem in TwinSight was converting dense CAD-derived geometry into assets suitable for real-time WebGL inspection.

The raw assets were not designed for browser-based rendering. They had to preserve enough detail to remain technically meaningful, but also be simplified enough to support interaction.

I approached it by analyzing the geometry, identifying what detail was necessary for recognition and assembly understanding, simplifying or rebuilding parts where needed, preparing assets in Blender and integrating them into Unity.

The result was an optimized model that preserved the inspection value of the assembly while becoming more appropriate for real-time visualization.
\end{mdcode}




\medskip\hrule\medskip




\section{40. “What would you do if you joined our team?”}




\subsection{Answer}




\begin{mdcode}
Bloque de codigo: text
First, I would understand the project constraints: target platform, performance budget, asset pipeline, current architecture, user needs and delivery priorities.

Then I would look for where I can contribute fastest: Unity/C# interaction systems, technical UI, asset optimization, visualization tools, documentation or workflow automation.

I would avoid making assumptions too early. My first goal would be to understand how the team currently builds and ships, then contribute in a way that improves the existing workflow rather than fighting it.
\end{mdcode}




\medskip\hrule\medskip




\section{41. “How do you learn new tools?”}




\subsection{Answer}




\begin{mdcode}
Bloque de codigo: text
I learn best by connecting tools to a concrete output.

For example, I do not treat a course as valuable just because I completed it. I try to convert learning into an artifact: a prototype, a breakdown, a before/after optimization, a demo or documentation.

That is the same logic I use for Rebelway, Blender, Unity or Python automation. The value is not the course itself; it is the public evidence it helps produce.
\end{mdcode}




\medskip\hrule\medskip




\section{42. Questions to ask interviewers}




\subsection{General}




\begin{mdcode}
Bloque de codigo: text
What kind of real-time 3D or Unity systems would I be working on first?
How is the team structured between developers, artists, technical artists and designers?
What platforms are most important for this role?
What are the main technical constraints right now?
How do you evaluate success for this role in the first 3 months?
\end{mdcode}




\subsection{Unity / Technical Art}




\begin{mdcode}
Bloque de codigo: text
What Unity version and render pipeline are you using?
How do you currently handle profiling and optimization?
How are technical artists involved in runtime systems?
How much of the role is tools, shaders, asset pipeline or implementation?
Do you expect this person to write production C#?
\end{mdcode}




\subsection{WebGL}




\begin{mdcode}
Bloque de codigo: text
Is WebGL a primary deployment target or secondary?
What are the current build-size and performance constraints?
How do you handle asset loading and compression?
What browsers/devices are supported?
Do you use Unity WebGL, Three.js, custom WebGL or another stack?
\end{mdcode}




\subsection{Technical visualization / digital twin}




\begin{mdcode}
Bloque de codigo: text
What kinds of technical assets or CAD data do you work with?
How do you process CAD or engineering data for real-time visualization?
Is the system connected to live data, IoT, APIs or static assets?
What level of simulation accuracy is expected?
What parts are visualization vs engineering calculation?
\end{mdcode}




\subsection{Contractor}




\begin{mdcode}
Bloque de codigo: text
Is this role structured as contractor, employee, B2B or EOR?
What is the expected weekly availability?
What time-zone overlap is required?
What are the payment terms and currency?
How is IP ownership handled?
Are there exclusivity or non-compete clauses?
\end{mdcode}




\medskip\hrule\medskip




\section{43. Questions not to ask too early}




Avoid in first interview unless they bring it up:




\begin{itemize}
\item vacation details;
\item exact benefits if contractor structure is unclear;
\item “how soon can I get promoted?”;
\item relocation before role fit is clear;
\item asking for visa sponsorship before explaining current contractor availability;
\item deep salary negotiation before understanding role scope.
\end{itemize}




\medskip\hrule\medskip




\section{44. Red flag responses during interviews}




Watch for:




\begin{center}
\scriptsize
\begin{longtable}{>{\raggedright\arraybackslash}p{0.455\linewidth} >{\raggedright\arraybackslash}p{0.455\linewidth}}
\toprule
Signal & Risk \\
\midrule
\endfirsthead
\toprule
Signal & Risk \\
\midrule
\endhead
“We pay after client approval” & Payment risk \\
“The test should take a weekend” & Free labor risk \\
“We need full-time but contractor only” & Misclassification risk \\
“No contract needed” & Legal/payment risk \\
“Equity instead of pay” & Cashflow risk \\
“We own all your side projects” & IP risk \\
“Remote, but must live in the US/EU” & Authorization blocker \\
“We need senior experience” & Misfit \\
“Can you do design, web, AI, 3D, marketing?” & Scope chaos \\
\bottomrule
\end{longtable}
\end{center}




\medskip\hrule\medskip




\section{45. Final preparation checklist before interviews}




\begin{itemize}
\item [ ] 30-second pitch memorized.
\item [ ] 60-second TwinSight explanation memorized.
\item [ ] AI-assisted development answer ready.
\item [ ] Low formal experience answer ready.
\item [ ] Remote Colombia answer ready.
\item [ ] EU authorization answer ready.
\item [ ] Salary expectation answer ready.
\item [ ] Portfolio link ready.
\item [ ] GitHub link ready.
\item [ ] Demo video ready.
\item [ ] Questions for interviewer ready.
\item [ ] Role-specific technical notes reviewed.
\item [ ] Contract red flags reviewed.
\end{itemize}




\medskip\hrule\medskip




\section{46. Next file dependency}




The next file should be:




\begin{mdcode}
Bloque de codigo: text
24_offer_evaluation_scorecard.md
\end{mdcode}




It should include:




\begin{itemize}
\item offer scoring model;
\item weighted criteria;
\item gross-to-net model;
\item contractor vs employee evaluation;
\item red-flag checklist;
\item acceptance / negotiation / rejection thresholds;
\item compensation scenarios;
\item EU mobility considerations;
\item Colombia compliance checklist.
\end{itemize}




No Deep Research or agent required.



